% Lösungen Einheit 3 von 4 – Einheit 3 (zusammenbau v0.1)
\einheitenkopf{Einheit 3 von 4 · Einheit 3}

\erg{23}{a) Binomialformel – Bernoulli-Kette mit fester Trefferwahrscheinlichkeit \quad b) $p = \frac{50}{100} = 0{,}5$; $\sigma = \sqrt{100 \cdot 0{,}5 \cdot 0{,}5} = 5$ \quad c) $12 \cdot (1 - p) = 3$, also $1 - p = 0{,}25$ und $p = 0{,}75$; $n = 12 : 0{,}75 = 16$ \quad d) $E(X) = 10 \cdot 0{,}1 = 1$; Summand $(10 - 1)^2 \cdot 0{,}1 = 8{,}1$, also Varianz mindestens 8{,}1 und $\sigma \ge \sqrt{8{,}1} \approx 2{,}85 > 2{,}5$ \quad e) $n$ steht unter der Wurzel, doppeltes $\sigma$ braucht vierfaches $n$: $n = 4 \cdot 50 = 200$ \quad f) Der Bogen reicht von $p = 0$ bis $p = 1$, also $1 : 20 = 0{,}05$ je Kästchen; bei $p = 0{,}25$ ist $\sigma = \sqrt{1\,200 \cdot 0{,}25 \cdot 0{,}75} = 15$, fünf Kästchen, also 3 je Kästchen; das Maximum $\sqrt{300} \approx 17{,}3$ liegt auf keiner Gitterlinie und taugt nicht zum Skalieren}
\erg{24}{a) $\sqrt{n \cdot 0{,}1 \cdot 0{,}9} = 6$, quadriert $0{,}09 \cdot n = 36$, also $n = 400$}
\erg{25}{a) $\mu_X = 2$, $p_X = 0{,}2$, Varianz $1{,}6$; $\mu_Y = 5$, $p_Y = 0{,}5$, Varianz $2{,}5$; Verhältnis $\frac{1{,}6}{2{,}5} = 0{,}64$}
\erg{26}{a) Fehler: Emil setzt $\sigma$ statt $\sigma^2$ ein; richtig $36 = 40 \cdot (1 - p)$, $1 - p = 0{,}9$, also $p = 0{,}1$ und $n = 40 : 0{,}1 = 400$ \quad b) $n \cdot p \cdot (1 - p) = np - np^2$ ist ein quadratischer Term in $p$ mit negativem Faktor vor $p^2$, also eine nach unten geöffnete Parabel; ihr Scheitel in der Mitte, bei $p$ gleich einhalb, ist die größte Streuung. \quad c) $\sigma$: 3; 4; 5; 4; 3 \quad d) $\mu = 500 \cdot 0{,}02 = 10$; $\sigma = \sqrt{500 \cdot 0{,}02 \cdot 0{,}98} = \sqrt{9{,}8} \approx 3{,}13$}
